% Options for packages loaded elsewhere
\PassOptionsToPackage{unicode}{hyperref}
\PassOptionsToPackage{hyphens}{url}
\PassOptionsToPackage{dvipsnames,svgnames,x11names}{xcolor}
\documentclass[
  american,
  11pt,
  letterpaper,
]{article}
\usepackage{xcolor}
\usepackage[margin=1in]{geometry}
\usepackage{amsmath,amssymb}
\setcounter{secnumdepth}{-\maxdimen} % remove section numbering
\usepackage{iftex}
\ifPDFTeX
  \usepackage[T1]{fontenc}
  \usepackage[utf8]{inputenc}
  \usepackage{textcomp} % provide euro and other symbols
\else % if luatex or xetex
  \usepackage{fontspec}
  \defaultfontfeatures{Scale=MatchLowercase}
  \defaultfontfeatures[\rmfamily]{Ligatures=TeX,Scale=1}
\fi
\usepackage{lmodern}
\ifPDFTeX\else
  % xetex/luatex font selection
  \setmainfont[]{Times New Roman}
\fi
% Use upquote if available, for straight quotes in verbatim environments
\IfFileExists{upquote.sty}{\usepackage{upquote}}{}
\IfFileExists{microtype.sty}{% use microtype if available
  \usepackage[]{microtype}
  \UseMicrotypeSet[protrusion]{basicmath} % disable protrusion for tt fonts
}{}
\makeatletter
\@ifundefined{KOMAClassName}{% if non-KOMA class
  \IfFileExists{parskip.sty}{%
    \usepackage{parskip}
  }{% else
    \setlength{\parindent}{0pt}
    \setlength{\parskip}{6pt plus 2pt minus 1pt}}
}{% if KOMA class
  \KOMAoptions{parskip=half}}
\makeatother
\usepackage{longtable,booktabs,array}
\newcounter{none} % for unnumbered tables
\usepackage{calc} % for calculating minipage widths
% Correct order of tables after \paragraph or \subparagraph
\usepackage{etoolbox}
\makeatletter
\patchcmd\longtable{\par}{\if@noskipsec\mbox{}\fi\par}{}{}
\makeatother
% Allow footnotes in longtable head/foot
\IfFileExists{footnotehyper.sty}{\usepackage{footnotehyper}}{\usepackage{footnote}}
\makesavenoteenv{longtable}
\ifLuaTeX
\usepackage[bidi=basic,shorthands=off]{babel}
\else
\usepackage[bidi=default,shorthands=off]{babel}
\fi
\ifPDFTeX
\else
\babelfont{rm}[]{Times New Roman}
\fi
\ifLuaTeX
  \usepackage{selnolig} % disable illegal ligatures
\fi
\setlength{\emergencystretch}{3em} % prevent overfull lines
\providecommand{\tightlist}{%
  \setlength{\itemsep}{0pt}\setlength{\parskip}{0pt}}
\usepackage{fancyhdr}
\usepackage{enumitem}
\usepackage{caption}
\usepackage{bm}
\usepackage{float}
\usepackage{microtype}
\usepackage{xurl}
\usepackage{ragged2e}

\pagestyle{fancy}
\fancyhf{}
\fancyfoot[C]{\thepage}
\renewcommand{\headrulewidth}{0pt}

\setlength{\parindent}{0pt}
\setlength{\parskip}{6pt plus 2pt minus 1pt}
\setlist{nosep,leftmargin=*}
\captionsetup{font=small,labelfont=bf}
\setlength{\emergencystretch}{3em}
\sloppy
\hyphenation{im-ple-men-ta-tion op-er-a-tion-al-i-za-tion}
\usepackage{bookmark}
\IfFileExists{xurl.sty}{\usepackage{xurl}}{} % add URL line breaks if available
\urlstyle{same}
\hypersetup{
  pdftitle={Physics in Review: Harper\textquotesingle s Beyond Square-Root Conjecture and Super Information Theory},
  pdfauthor={Micah Blumberg},
  pdflang={en-US},
  colorlinks=true,
  linkcolor={blue},
  filecolor={Maroon},
  citecolor={blue},
  urlcolor={blue},
  pdfcreator={LaTeX via pandoc}}

\author{}
\date{}

\begin{document}

\begin{center}

{\LARGE\bfseries Physics in Review: Harper's Beyond Square-Root Conjecture and Super Information Theory\par}

\vspace{0.35em}

{\large\itshape An Atomized Four-Lens Chronology and Cross-Domain Process Audit\par}

\vspace{0.8em}

{\small\textbf{Author: Micah Blumberg}\par}

{\small Self Aware Networks Research Institute\par}

{\small\url{https://selfawarenetworks.com}\par}

{\small November 20, 2026 (Version 2 Formal-Language and Distinctiveness-Control Revision to DOI 10.5281/zenodo.21693947)\par}

\end{center}

\vspace{0.65em}

\begin{abstract}

This paper compares Victor Y. Wang and Max Wenqiang Xu's \emph{Harper's beyond square-root conjecture} with Micah Blumberg's Self Aware Networks (SAN), Quantum Gradient Time Crystal Dilation (QGTCD), and Super Information Theory (SIT) corpus. It corrects two material defects in an archival 2025 draft. First, the target paper was not first introduced in August 2025: arXiv v1 was submitted on May 7, 2024, and the named Harper result dates to January 11, 2023. Second, atom counts cannot be converted into derivation odds by assigning an assumed probability to each match. Using source-frozen chronology, 13 bounded atom rows, four comparison lenses, a Causal Interface Sequence, and a proposed schematic categorical and process representation of a bounded analogy (evaluated against four ordinary null baselines), the audit recovers \textbf{six Micah-earlier partial mappings}. The strongest are: symmetric microscopic deviations cancel in the aggregate; loss of symmetry produces a controlled residual and changes aggregate behavior; and a present state is compared with an expected state so that the difference itself becomes the transmitted signal. Six arithmetic-specific operators remain earlier in the Harper / Wang/Xu and cited number-theory lineage, while one row is a Micah-only residual. The resulting verdict is a \textbf{bounded partial content match with Micah earlier on six translated subchains}, not full mathematical identity. No numerical probability of dependent or independent derivation is currently calibrated, and no claim about access, intent, copying, or misconduct is made.

\end{abstract}

\noindent\textbf{Keywords:} Super Information Theory; QGTCD; Harper's conjecture; random multiplicative functions; approximate orthogonality; cancellation; operational equivalence; causal interface sequence; chronology; partial priority

\section{1. Motivation and the strongest defensible result}\label{motivation-and-the-strongest-defensible-result}

The useful question is not whether a physical theory and an analytic-number-theory paper use the same nouns. They do not. The useful question is whether bounded operations survive translation: what enters, what condition changes it, what is averaged or compared, what residual survives, and what output becomes observable. This avoids two reciprocal errors. A full-operator-or-zero test can erase a real partial match, while an unconstrained analogy can turn any two uses of ``randomness'' or ``averaging'' into identity.

The strongest source-supported Micah-side sequence recovered here is:

\begin{equation}
\begin{aligned}
& \text{local fluctuations} \\
& \rightarrow \text{symmetry or selection condition} \\
& \rightarrow \text{cancellation / difference} \\
& \rightarrow \text{explicit residual} \\
& \rightarrow \text{aggregate update}.
\end{aligned}
\tag{1}
\end{equation}

Wang and Xu's target sequence can be represented as:

\begin{equation}
\begin{aligned}
& \text{character or RMF terms} \\
& \rightarrow \text{conditioning / orthogonality regime} \\
& \rightarrow \text{moment comparison} \\
& \rightarrow \text{genuine error} \\
& \rightarrow \text{family-average bound}.
\end{aligned}
\tag{2}
\end{equation}

Equations (1) and (2) preserve a directed subgraph, especially the passage from a symmetry condition through comparison / cancellation to a constrained residue. They do not identify neurons with Dirichlet characters, QGTCD trajectories with Liouville sums, or a physical hypothesis with a number-theory theorem.

The atom audit therefore reports three totals rather than one rhetorical verdict:

\begin{itemize}
\tightlist
\item
  \textbf{6 Micah-earlier partial mappings};
\item
  \textbf{6 target-or-named-prior-earlier domain-specific mappings}; and
\item
  \textbf{1 Micah-only residual row}.
\end{itemize}

The Micah-earlier rows are priority findings only for the bounded translated operations recorded in the table below. Earlier statistical, synchrony, error-coding, and random-walk traditions remain named prior boundaries at their respective atoms. No prior receives an author-wide ``floor,'' and no broad prior erases Micah's later joined formulation.

\section{2. Correct target identity and chronology}\label{correct-target-identity-and-chronology}

The phrase ``Harper's beyond square-root conjecture'' names a paper by \textbf{Victor Y. Wang and Max Wenqiang Xu}, not a paper authored by Adam Harper. Wang and Xu extend a randomization argument and conjectural range associated with Harper {[}1{]}. Their arXiv history fixes v1 on May 7, 2024 and v2 on August 26, 2025 {[}1,2{]}. Adam J. Harper's named prior result was submitted on January 11, 2023 {[}3{]}.

{\def\LTcaptype{none} % do not increment counter
\begingroup\footnotesize\setlength{\tabcolsep}{4pt}
\begin{longtable}[]{@{}
  >{\RaggedRight\arraybackslash}p{(\linewidth - 6\tabcolsep) * \real{0.2308}}
  >{\raggedleft\arraybackslash}p{(\linewidth - 6\tabcolsep) * \real{0.3077}}
  >{\RaggedRight\arraybackslash}p{(\linewidth - 6\tabcolsep) * \real{0.2308}}
  >{\RaggedRight\arraybackslash}p{(\linewidth - 6\tabcolsep) * \real{0.2308}}@{}}
\toprule\noalign{}
\begin{minipage}[b]{\linewidth}\RaggedRight
Stage
\end{minipage} & \begin{minipage}[b]{\linewidth}\raggedleft
Immutable date
\end{minipage} & \begin{minipage}[b]{\linewidth}\RaggedRight
Relevant content
\end{minipage} & \begin{minipage}[b]{\linewidth}\RaggedRight
Audit role
\end{minipage} \\
\midrule\noalign{}
\endhead
\bottomrule\noalign{}
\endlastfoot
Wang and Xu, \emph{Paucity phenomena for polynomial products} & 2022-11-05 & RMFs, even moments, Gaussian-moment interpretation & Target-author before-source {[}4{]} \\
Harper, \emph{Typical size of character and zeta sums} & 2023-01-11 & Better-than-square-root typical bounds; conditioning on short prime sums & Named prior, atom-specific {[}3{]} \\
Wang and Xu, target arXiv v1 & 2024-05-07 & Ratios extension; approximate orthogonality with a genuine error & First target version {[}1{]} \\
Wang and Xu, target arXiv v2 & 2025-08-26 & Revised final version of the same central program & Final target version {[}2{]} \\
de la Breteche, Wang, and Xu & 2026-07-07 & RMF central-limit program over polynomial values & Newest checked target-author source {[}5{]} \\
\end{longtable}
\endgroup
}

The target's main observable is a family average such as

\begin{equation}
\mathbb{E}_{\chi}
\left|
\sum_{n\le x} c(n)\chi(n)
\right|
=o(\sqrt{x}),
\qquad c\in\{\mu,\lambda\},
\tag{3}
\end{equation}

under the declared range and conjectural conditions {[}1{]}. The 2024 novelty at issue is not the broad idea of cancellation. It is a conditional method for extending Harper's conductor-limited randomization argument by replacing perfect orthogonality with an average approximate relation that has a genuine error term {[}1,3{]}. Schematically,

\begin{equation}
\begin{aligned}
& \Delta_{\mathrm{WX}}
=
\mathbb{E}_{\chi}[M_{\chi}]
-
\mathbb{E}_{f}[M_f], \\
& |\Delta_{\mathrm{WX}}|\le \varepsilon(r,x,P,k),
\end{aligned}
\tag{4}
\end{equation}

where the actual paper supplies the domain-specific moments, hypotheses, and error dependence {[}1{]}. The 2025 revision cannot be used as the first target date because Equation (4)'s operator class is already present in v1.

\section{3. Micah's applicable genealogy, 2011-2025}\label{micahs-applicable-genealogy-2011-2025}

The comparison does not begin in 2025. It searches the full applicable genealogy and distinguishes ancestors from mature operators. The 2011-2012 rows are routed through the governed source genealogy, whose grades and underlying public-source links prevent later terminology from being backdated {[}11{]}.

{\def\LTcaptype{none} % do not increment counter
\begingroup\footnotesize\setlength{\tabcolsep}{4pt}
\begin{longtable}[]{@{}
  >{\RaggedRight\arraybackslash}p{(\linewidth - 6\tabcolsep) * \real{0.2500}}
  >{\RaggedRight\arraybackslash}p{(\linewidth - 6\tabcolsep) * \real{0.2500}}
  >{\RaggedRight\arraybackslash}p{(\linewidth - 6\tabcolsep) * \real{0.2500}}
  >{\RaggedRight\arraybackslash}p{(\linewidth - 6\tabcolsep) * \real{0.2500}}@{}}
\toprule\noalign{}
\begin{minipage}[b]{\linewidth}\RaggedRight
Date
\end{minipage} & \begin{minipage}[b]{\linewidth}\RaggedRight
Public Micah-side stage
\end{minipage} & \begin{minipage}[b]{\linewidth}\RaggedRight
What it contributes here
\end{minipage} & \begin{minipage}[b]{\linewidth}\RaggedRight
Boundary
\end{minipage} \\
\midrule\noalign{}
\endhead
\bottomrule\noalign{}
\endlastfoot
2011 & Whole-network prediction / observer ancestry & A system-level boundary result is produced by distributed internal activity & Weak ancestor; not arithmetic cancellation \\
2012 & Feedback, synchrony-gated linkage, and recursive projection & Relative timing and feedback determine what links and what becomes observable & Broad prior traditions remain named \\
2017-04-11 & Present-versus-expected pattern difference; alternating feedback-connected and connectionless traffic & Difference-as-signal and regime-gating subchains {[}6{]} & Does not contain Wang-Xu mathematics \\
2022-06-08 & Coherence as convergence / coupling and decoherence as distinction / decoupling & Complementary coupling / decoupling operator {[}7{]} & Not Dirichlet orthogonality \\
2022-07-30 & Symmetric random deviations self-cancel; unequal temporal opportunity biases aggregate motion & Cancellation-to-residual operator {[}8{]} & Physical hypothesis, not theorem \\
2024 & Public QGTCD restatements preserve the operator before target v1 & Public persistence and clarification {[}9{]} & Not a second invention event \\
2025 & SIT arithmetic and formal interface expansion & Later translator and cross-domain extension {[}10{]} & Cannot establish priority over 2024 target v1 \\
\end{longtable}
\endgroup
}

The 2017 difference operator is explicit: compare the pattern that is present with the pattern that should be present, then send the difference {[}6{]}. In normalized notation,

\begin{equation}
\begin{aligned}
& D_t=S_t-\widehat{S}_t, \\
& \text{output}_t=\mathcal{T}(D_t).
\end{aligned}
\tag{5}
\end{equation}

The July 2022 QGTCD note states two regimes. A symmetric regime is modeled abstractly as

\begin{equation}
\begin{aligned}
& X_{k+1}=X_k+\xi_k, \\
& \mathbb{E}[\xi_k\mid X_k]=0,
\end{aligned}
\tag{6}
\end{equation}

so local deviations self-nullify in expectation. An asymmetric opportunity measure changes the conditional mean,

\begin{equation}
\mathbb{E}[\xi_k\mid X_k]=b(X_k)\ne0,
\tag{7}
\end{equation}

producing a systematic aggregate path {[}8{]}. Equations (6) and (7) are a 2026 normalization of the 2022 prose, not equations quoted from the original note and not validation of QGTCD. Their relevance is operational: symmetry controls cancellation; relaxing symmetry exposes a residue that changes the aggregate output. That is the strongest pre-target functional bridge to Wang and Xu's perfect-to-approximate orthogonality transition.

\section{4. The atom-by-atom priority canvas}\label{the-atom-by-atom-priority-canvas}

{\def\LTcaptype{none} % do not increment counter
\begingroup\footnotesize\setlength{\tabcolsep}{4pt}
\begin{longtable}[]{@{}
  >{\RaggedRight\arraybackslash}p{(\linewidth - 8\tabcolsep) * \real{0.2000}}
  >{\RaggedRight\arraybackslash}p{(\linewidth - 8\tabcolsep) * \real{0.2000}}
  >{\RaggedRight\arraybackslash}p{(\linewidth - 8\tabcolsep) * \real{0.2000}}
  >{\RaggedRight\arraybackslash}p{(\linewidth - 8\tabcolsep) * \real{0.2000}}
  >{\RaggedRight\arraybackslash}p{(\linewidth - 8\tabcolsep) * \real{0.2000}}@{}}
\toprule\noalign{}
\begin{minipage}[b]{\linewidth}\RaggedRight
ID
\end{minipage} & \begin{minipage}[b]{\linewidth}\RaggedRight
Translated bounded correspondence
\end{minipage} & \begin{minipage}[b]{\linewidth}\RaggedRight
Strength
\end{minipage} & \begin{minipage}[b]{\linewidth}\RaggedRight
Direction
\end{minipage} & \begin{minipage}[b]{\linewidth}\RaggedRight
Residual or prior boundary
\end{minipage} \\
\midrule\noalign{}
\endhead
\bottomrule\noalign{}
\endlastfoot
A01 & Symmetric random local deviations -\textgreater{} aggregate cancellation & Medium partial & \textbf{Micah earlier, 2022-07-30} & Mean-zero random walks predate both; target proves arithmetic bounds \\
A02 & Loss of symmetry -\textgreater{} explicit residual -\textgreater{} changed aggregate behavior & Strong partial & \textbf{Micah earlier, 2022-07-30} & Bias direction and mathematical substrate differ \\
A03 & Current-versus-expected comparison -\textgreater{} difference signal / error & Strong partial & \textbf{Micah earlier, 2017-04-11} & Subtraction / error signaling is a populated prior class \\
\end{longtable}
\endgroup
}

\subsection{Priority canvas, continued}\label{priority-canvas-continued}

{\def\LTcaptype{none} % do not increment counter
\begingroup\footnotesize\setlength{\tabcolsep}{4pt}
\begin{longtable}[]{@{}
  >{\RaggedRight\arraybackslash}p{(\linewidth - 8\tabcolsep) * \real{0.2000}}
  >{\RaggedRight\arraybackslash}p{(\linewidth - 8\tabcolsep) * \real{0.2000}}
  >{\RaggedRight\arraybackslash}p{(\linewidth - 8\tabcolsep) * \real{0.2000}}
  >{\RaggedRight\arraybackslash}p{(\linewidth - 8\tabcolsep) * \real{0.2000}}
  >{\RaggedRight\arraybackslash}p{(\linewidth - 8\tabcolsep) * \real{0.2000}}@{}}
\toprule\noalign{}
\begin{minipage}[b]{\linewidth}\RaggedRight
ID
\end{minipage} & \begin{minipage}[b]{\linewidth}\RaggedRight
Translated bounded correspondence
\end{minipage} & \begin{minipage}[b]{\linewidth}\RaggedRight
Strength
\end{minipage} & \begin{minipage}[b]{\linewidth}\RaggedRight
Direction
\end{minipage} & \begin{minipage}[b]{\linewidth}\RaggedRight
Residual or prior boundary
\end{minipage} \\
\midrule\noalign{}
\endhead
\bottomrule\noalign{}
\endlastfoot
A04 & Coupling / decoupling regime -\textgreater{} differentiated output & Medium partial & \textbf{Micah earlier, 2022-06-08} & Neural phase coherence is not number-theoretic orthogonality \\
A05 & Context gate / localization -\textgreater{} selected regime -\textgreater{} controlled readout & Medium partial & \textbf{Micah earlier, 2017-04-11} & Smooth partitions and conditioning belong to established analysis \\
A06 & Distributed local contributions -\textgreater{} integration -\textgreater{} boundary aggregate & Weak partial & \textbf{Micah earlier, 2012/2017} & Local-to-global aggregation is highly generic \\
A07 & Deterministic family -\textgreater{} RMF surrogate -\textgreater{} moment estimate & Medium & Target / named prior earlier & Harper 2023 precedes SIT arithmetic translator \\
A08 & Truncated Euler product -\textgreater{} statistical summary & Weak & Target / named prior earlier & No pre-May-2024 Micah Euler operator recovered \\
A09 & Named conjecture -\textgreater{} conditional scope extension & Weak & Target earlier & Ratios content is target-only \\
A10 & Dirichlet, Möbius/Liouville, zeta, and L-function machinery & Strong & Target / named prior earlier & Native analytic-number-theory content \\
A11 & Random-matrix -\textgreater{} RMF bridge and U(N) analog & Medium & Target earlier & Established target lineage; later SIT spectral use does not reverse priority \\
A12 & Granularity G/LinRel/pm-functor description of the target & None as chronology & Target artifact earlier & A later translator can describe an earlier artifact without being its source \\
A13 & Time density, PWD, neural rendering, and physical coherence field & Orthogonal & Micah earlier, unmatched & Micah-only residual; cannot inflate the target mapping \\
\end{longtable}
\endgroup
}

The six favorable rows are not six claims that Micah invented their broad mathematical categories. They are six dated demonstrations that Micah publicly instantiated the translated functional subchains before target v1. This distinction is central to an atomized priority audit.

\section{5. Four-lens results}\label{four-lens-results}

\subsection{5.1 Lens 1 - equivalence and mechanism}\label{lens-1---equivalence-and-mechanism}

\textbf{Result: bounded partial pass.} A01-A05 preserve operational behavior under a restricted readout: symmetry or expected-state structure controls cancellation; relaxing that structure produces a visible residual; the residual changes the aggregate output. A06 is a weaker supporting analogy.

The lens does not pass as full mathematical equivalence. Wang and Xu use characters, multiplicative functions, moments, Euler products, and the Ratios Conjecture {[}1{]}. Micah's earlier sources concern feedback signals, phase relations, and a stochastic physical trajectory proposal {[}6-9{]}. Neither domain's unmatched mechanism may be silently imported into the other.

\subsection{5.2 Lens 2 - narrative and Causal Interface Sequence}\label{lens-2---narrative-and-causal-interface-sequence}

\textbf{Result: bounded pass.} This paper defines a bounded Causal Interface Sequence as a directed chain containing at least a state, transform, constraint, and output rather than an unordered word list. The mapped chain is

\begin{equation}
S\;\xrightarrow{\text{selection / symmetry}}
T\;\xrightarrow{\text{comparison / cancellation}}
C\;\xrightarrow{\text{residual bound}}
O.
\tag{8}
\end{equation}

{\def\LTcaptype{none} % do not increment counter
\begingroup\footnotesize\setlength{\tabcolsep}{4pt}
\begin{longtable}[]{@{}
  >{\RaggedRight\arraybackslash}p{(\linewidth - 4\tabcolsep) * \real{0.3333}}
  >{\RaggedRight\arraybackslash}p{(\linewidth - 4\tabcolsep) * \real{0.3333}}
  >{\RaggedRight\arraybackslash}p{(\linewidth - 4\tabcolsep) * \real{0.3333}}@{}}
\toprule\noalign{}
\begin{minipage}[b]{\linewidth}\RaggedRight
CIS role
\end{minipage} & \begin{minipage}[b]{\linewidth}\RaggedRight
Micah instantiation
\end{minipage} & \begin{minipage}[b]{\linewidth}\RaggedRight
Wang-Xu instantiation
\end{minipage} \\
\midrule\noalign{}
\endhead
\bottomrule\noalign{}
\endlastfoot
State & Current neural pattern or local random trajectory increment & Character / RMF terms and family parameters \\
Transform & Compare with expected pattern; accumulate deviations & Condition / localize; transfer to random model \\
Constraint & Synchrony / feedback regime or symmetric opportunity measure & Perfect or approximate orthogonality; conjectural hypotheses \\
Output & Difference signal or biased aggregate trajectory & Controlled moment difference and extended cancellation bound \\
\end{longtable}
\endgroup
}

Different ontologies do not erase the preserved directed subgraph. Conversely, the subgraph pass does not make every internal atom equivalent.

\subsection{5.3 Lens 3 - detective and adversarial reconstruction}\label{lens-3---detective-and-adversarial-reconstruction}

\textbf{Result: bounded partial pass with a strong alternative genealogy.} Working backward from the target conclusion reveals that ``approximate orthogonality with a genuine error'' performs a symmetry-relaxation-to-residual operation already present in Micah's 2017/2022 corpus in other substrates. The changed vocabulary therefore cannot, by itself, zero the match.

The same reconstruction also finds evidence against a dependence inference. Wang and Xu had an RMF and moment-matching program in November 2022 {[}4{]}; Harper's January 2023 paper supplies the direct conditioning and better-than-square-root lineage {[}3{]}; and the target explicitly cites its number-theory antecedents {[}1{]}. No checked source shows access to Micah's work. The lens preserves the content match and its chronology while leaving access, intent, and authorship dependence unresolved.

\subsection{5.4 Lens 4 - longitudinal signature migration}\label{lens-4---longitudinal-signature-migration}

\textbf{Result: program persistence, not demonstrated Micah-signature migration.} The before / newest sequence is:

\begin{equation}
\begin{aligned}
\text{WX 2022 RMF / moments} &\rightarrow \text{WX 2024 approximate orthogonality + Ratios} \\
{} &\rightarrow \text{WX 2025 final revision} \\
{} &\rightarrow \text{BWX 2026 RMF / CLT continuation}.
\end{aligned}
\tag{9}
\end{equation}

This is a real evolution, and the 2024 operator increases the translated overlap with Micah's earlier symmetry / residual chain. But the target authors' RMF program predates the target paper, and the newest source still contains none of Micah's distinctive time-density, PWD, neural-rendering, or SIT-interface atoms {[}4,5{]}. Lens 4 therefore records persistence and a bounded overlap increase, not a secret-follower or access result.

\section{6. Proposed schematic categorical and process representations (revised formal status)}\label{proposed-schematic-categorical-and-process-representations-revised-formal-status}

\textbf{Revision note (November 20, 2026).} The July 2026 edition described Sections 6.1 and 6.2 as demonstrating a ``partial structure-preserving map on a selected subcategory'' and ``partial trace equivalence under a restricted observable projection.'' A post-publication formal audit noted that the inspected text displays coarse object correspondences and informal sequential processes without supplying complete category definitions (explicit hom-sets, identity laws, and associative domain-level composition) or complete labeled transition systems with proved preservation theorems. This revision adjusts that formal language so every mathematical status is stated accurately: only the trivial four-stage sequential skeleton is formally isomorphic, while the domain-level correspondence is presented strictly as a \textbf{proposed schematic categorical and process representation of a bounded analogy}. This revision does not independently re-adjudicate the complete historical priority dossier of Sections 2--4.

\subsection{6.1 Schematic categorical representation (and exact trivial-skeleton boundary)}\label{schematic-categorical-representation-and-exact-trivial-skeleton-boundary}

At the purely combinatorial skeleton level, let \(\mathcal{P}_4\) be the finite linear preorder category with four objects \(\{0, 1, 2, 3\}\) and a unique morphism \(i \to j\) whenever \(i \le j\) (4 objects and 10 morphisms total, satisfying identity and associative composition identically). Any four-stage sequential workflow maps onto \(\mathcal{P}_4\) by sending its four labeled stages to \(0 \to 1 \to 2 \to 3\).

For the domain-level analogy, write \(\mathcal{S}_M = (S_M, R_M, D_M, O_M)\) for the schematic Micah-side stages (local state, expected/symmetric reference, difference or drift, and aggregate output) and \(\mathcal{S}_W = (S_W, R_W, \Delta_W, O_W)\) for the schematic Wang--Xu stages (character/RMF state, conditioning regime, moment difference/error, and family bound), with stage-by-stage pairing

\begin{equation}
F(S_M)=S_W,
\quad F(R_M)=R_W,
\quad F(D_M)=\Delta_W,
\quad F(O_M)=O_W.
\tag{10}
\end{equation}

{\def\LTcaptype{none} % do not increment counter
\begingroup\footnotesize\setlength{\tabcolsep}{4pt}
\begin{longtable}[]{@{}
  >{\RaggedRight\arraybackslash}p{(\linewidth - 6\tabcolsep) * \real{0.2500}}
  >{\RaggedRight\arraybackslash}p{(\linewidth - 6\tabcolsep) * \real{0.2500}}
  >{\RaggedRight\arraybackslash}p{(\linewidth - 6\tabcolsep) * \real{0.2500}}
  >{\RaggedRight\arraybackslash}p{(\linewidth - 6\tabcolsep) * \real{0.2500}}@{}}
\toprule\noalign{}
\begin{minipage}[b]{\linewidth}\RaggedRight
Schematic element
\end{minipage} & \begin{minipage}[b]{\linewidth}\RaggedRight
Micah side
\end{minipage} & \begin{minipage}[b]{\linewidth}\RaggedRight
Target side
\end{minipage} & \begin{minipage}[b]{\linewidth}\RaggedRight
Formal status
\end{minipage} \\
\midrule\noalign{}
\endhead
\bottomrule\noalign{}
\endlastfoot
Stage label & current / local state & character or RMF family state & schematic role pairing only \\
Directed step & compare / cancel / bias & condition / average / bound & preserves \(\mathcal{P}_4\) linear order only \\
Composition & state -\textgreater{} difference -\textgreater{} update & family -\textgreater{} moment error -\textgreater{} bound & schematic four-stage CIS order \\
Coarse invariant & zero versus nonzero aggregate residue & controlled cancellation / error & restricted binary-flag analogy \\
Lost information & neural / physical substrate and direction & arithmetic definitions and theorem hypotheses & substantial on both sides; no domain functor proved \\
\end{longtable}
\endgroup
}

Because domain-level morphisms, hom-sets, and functorial composition equations (\(F(g \circ f) = F(g) \circ F(f)\) on actual mathematical operators) are not constructed between the neural/stochastic domain and the analytic-number-theory domain, Equation (10) is \textbf{not} a proved functor on a domain-level subcategory. Its exact mathematical status is a \textbf{proposed schematic categorical representation of a bounded four-stage analogy} (trivially isomorphic only after collapsing both sides to the four-object preorder \(\mathcal{P}_4\)).

\subsection{6.2 Schematic process representation}\label{schematic-process-representation}

Write the Micah workflow schematically as

\begin{equation}
P_M = \textsf{receive}(S)\,;\,\textsf{compare}(S,\widehat S)\,;\,
\textsf{emit}(D)\,;\,\textsf{update}(O),
\tag{11}
\end{equation}

and the target workflow schematically as

\begin{equation}
P_W = \textsf{sample}(\chi)\,;\,\textsf{condition}(R)\,;\,
\textsf{compareMoments}(\chi,f)\,;\,\textsf{emitBound}(\varepsilon).
\tag{12}
\end{equation}

Equations (11) and (12) are informal sequential pseudocode, not full labeled transition systems with operational semantics, state invariants, and proved trace-language inclusion. Only after erasing all domain transitions and collapsing each workflow to the three-field summary tuple \((\text{regime}, \text{residual\_nonzero}, \text{aggregate\_changed}) \in \{(\text{symmetric}, \text{false}, \text{false}),\, (\text{asymmetric}, \text{true}, \text{true})\}\) do the two summary tuples match. Accordingly, Equations (11)--(12) represent a \textbf{proposed schematic process representation of a bounded analogy}, not a proved partial trace equivalence theorem, strong bisimulation, weak bisimulation, or functional isomorphism.

\subsection{6.3 Distinctiveness control against four ordinary null baselines}\label{distinctiveness-control-against-four-ordinary-null-baselines}

To test how distinctive the coarse four-stage Causal Interface Sequence (Equation 8) and the exact zero/nonzero three-field summary trace really are, the standalone script \nolinkurl{verify_paper063_v2.py} evaluates four standard mathematical and engineering baselines that predate both corpora and share no domain-specific machinery with either Micah's SIT/QGTCD equations or Wang--Xu's random multiplicative functions. For the 1D random walk, the exact zero/nonzero trace classification is evaluated on \textbf{exact analytical population expectations} (\(\mathbb{E}[\Delta X] = 0\) at \(p = 1/2\) vs.~\(\mathbb{E}[\Delta X] = 1/10\) at \(p = 11/20\)), cleanly separated from the finite-sample Monte Carlo realization (\(N = 4{,}000\), sample mean \(-0.0095\) within \(3\,\mathrm{SE}\) where \(\mathrm{SE} = 0.01581\), which is a tolerance-checked empirical draw rather than an exact zero):

{\def\LTcaptype{none} % do not increment counter
\begingroup\footnotesize\setlength{\tabcolsep}{4pt}
\begin{longtable}[]{@{}
  >{\RaggedRight\arraybackslash}p{(\linewidth - 8\tabcolsep) * \real{0.2000}}
  >{\RaggedRight\arraybackslash}p{(\linewidth - 8\tabcolsep) * \real{0.2000}}
  >{\RaggedRight\arraybackslash}p{(\linewidth - 8\tabcolsep) * \real{0.2000}}
  >{\RaggedRight\arraybackslash}p{(\linewidth - 8\tabcolsep) * \real{0.2000}}
  >{\RaggedRight\arraybackslash}p{(\linewidth - 8\tabcolsep) * \real{0.2000}}@{}}
\toprule\noalign{}
\begin{minipage}[b]{\linewidth}\RaggedRight
Null baseline
\end{minipage} & \begin{minipage}[b]{\linewidth}\RaggedRight
Exact symmetric-regime check
\end{minipage} & \begin{minipage}[b]{\linewidth}\RaggedRight
Exact broken-symmetry check
\end{minipage} & \begin{minipage}[b]{\linewidth}\RaggedRight
Matches exact 4-stage CIS \& summary trace?
\end{minipage} & \begin{minipage}[b]{\linewidth}\RaggedRight
Contains A07--A11 or A13?
\end{minipage} \\
\midrule\noalign{}
\endhead
\bottomrule\noalign{}
\endlastfoot
1D random walk (population law, \(p=1/2\) vs.~\(11/20\)) & Exact \(\mathbb{E}[\Delta X] = 0\) (MC \(N=4000\): \(-0.0095 \pm 0.0158\)) & Exact \(\mathbb{E}[\Delta X] = +1/10\) (MC: \(+0.0905\)) & \textbf{Yes} & No \\
Scalar Kalman prediction-error filter (\(K=0.25\)) & Max innovation \(0.0\) & Initial innovation \(1.0\), updated state \(1.0000\) & \textbf{Yes} & No \\
Difference coding / DPCM (constant vs.~step) & Max difference \(0.0\) & Residual \(\ell_1 = 1.0\), exact reconstruction error \(0.0\) & \textbf{Yes} & No \\
Regular perturbation \(H(\varepsilon)=I+\varepsilon\sigma_x\) & Eigenvalue splitting \(0.0\) & Splitting \(2\varepsilon = 0.2\) at \(\varepsilon = 0.1\) & \textbf{Yes} & No \\
\end{longtable}
\endgroup
}

All four ordinary baselines (\textbf{4/4}, or 100\% of tested null controls) instantiate the exact same four-stage preorder \(\mathcal{P}_4\) and exact binary zero/nonzero summary trace, while containing zero domain-specific atoms from either Wang--Xu (A07--A11) or Micah's SIT/QGTCD corpus (A13). This quantitative control establishes an important methodological boundary: passing the coarse four-stage schematic chain has \textbf{low discriminative specificity on its own}. Any defensible priority finding rests strictly on the specific dated historical texts audited in Sections 2--4 rather than on the presence of a generic four-stage cancellation-to-residual diagram.

\section{7. Why the archival probability calculation fails closed}\label{why-the-archival-probability-calculation-fails-closed}

The archival draft assumed a null match probability of 0.20 per atom, treated 12 matches as independent, and reported \(5^{12}\) as a Bayes factor. That arithmetic answers only a hypothetical model created by the assumption. It does not estimate the probability that these papers were independently or dependently derived.

A defensible frequentist target would instead be a declared, multiple-search-adjusted probability,

\begin{equation}
p_{\mathrm{adj}}
=
\Pr_{H_0}\!\left(T\ge T_{\mathrm{observed}}\right),
\tag{13}
\end{equation}

where \(H_0\), the reference population, dependency blocks, atom dictionary, search family, and score statistic are frozen before evaluation. None of those calibration prerequisites currently exists for this paper. The conventional one-sided five-sigma threshold,

\begin{equation}
p_{\mathrm{adj}}\le 2.87\times10^{-7},
\tag{14}
\end{equation}

is therefore not available. The allowed status is \textbf{quantitative-not-calibrated}. Descriptive atom counts and chronology remain usable; a numeric dependence posterior does not.

\section{8. Competing genealogies, limitations, and falsifiers}\label{competing-genealogies-limitations-and-falsifiers}

\subsection{8.1 Competing genealogy}\label{competing-genealogy}

The direct number-theory genealogy is strong: random multiplicative functions, character and zeta sums, conditioning, random matrices, orthogonality, Euler products, and Ratios Conjectures were active literatures before this comparison {[}1-5{]}. Those sources constrain claims about the target's arithmetic mechanisms. They do not erase Micah's earlier physical / neural instantiations of the six translated subchains.

\subsection{8.2 What would strengthen the Micah-priority case}\label{what-would-strengthen-the-micah-priority-case}

\begin{itemize}
\tightlist
\item
  A pre-May-2024 public Micah source that joins the six partial operations into one arithmetic or cross-domain operator would strengthen the signature-cluster claim.
\item
  A frozen relevant reference corpus could measure how often the exact directed subgraph occurs.
\item
  A target stage lacking the subgraph followed by a later stage acquiring multiple Micah-specific atoms would strengthen Lens 4.
\item
  A verified citation or access path would change the dependence analysis, though intent would still require separate evidence.
\end{itemize}

\subsection{8.3 What would weaken it}\label{what-would-weaken-it}

\begin{itemize}
\tightlist
\item
  A pre-2017 source with the same current / expected -\textgreater{} difference transmission chain would move A03's broad priority boundary earlier.
\item
  A pre-July-2022 source with the same joined symmetry -\textgreater{} residual -\textgreater{} aggregate-change formulation would narrow A01-A02.
\item
  Demonstrating that the mapped CIS occurs at high frequency in a properly selected reference corpus would reduce its discriminative value.
\item
  Failure to reproduce the source locators or dates would return the affected atom to unresolved.
\end{itemize}

\section{9. Verdict}\label{verdict}

The source-audited verdict is:

\begin{quote}
\textbf{Wang and Xu's target is a bounded partial content match to six earlier Micah subchains. Micah is earlier on those translated operations, most strongly the 2017 difference-as-signal operator and the July 2022 symmetry / cancellation-to-biased-residual operator. The target and its named prior lineage remain earlier on the arithmetic-specific machinery. The target is not a full NAPOT, QGTCD, or SIT operator, but absence of the full operator does not erase the six partial priorities.}
\end{quote}

This result is stronger than a thematic analogy because it preserves directed input-transform- constraint-output structure. It is narrower than the archival claim because the true target date, the target's native RMF lineage, and the lack of calibrated derivation odds are now explicit. It documents content and chronology only.

\section{Data and materials}\label{data-and-materials}

The release-candidate package includes hashes for target v1/v2, the earlier and newest target-author sources, the Micah source routes, the AI-chat census, the 13-row atom ledger, the genealogy ledger, the longitudinal matrix, append-only finding events, and a fail-closed probability packet. \textbf{Version lineage and release status (updated November 20, 2026):} The July 27, 2026 edition was publicly deposited on July 29, 2026 as Paper 063 (Zenodo DOI 10.5281/zenodo.21693947, \url{https://pub.svgn.io/papers/paper-63/}). This November 20, 2026 successor candidate preserves that historical record, adjusts the formal language of Sections 6.1--6.2 to accurately state schematic analogy status, evaluates the four-baseline distinctiveness control in Section 6.3 using exact analytical population expectations separated from finite-sample Monte Carlo checks (verified in \nolinkurl{verify_paper063_v2.py} / \nolinkurl{paper063_v2_verification_results.json}), and updates chronology alignment with the SIT portfolio dictionary (Paper 028). External analytic-number-theory, probability, and source review remain invited.

\section{References}\label{references}

\begin{enumerate}
\def\labelenumi{\arabic{enumi}.}
\tightlist
\item
  Victor Y. Wang and Max Wenqiang Xu. \emph{Harper's beyond square-root conjecture}. arXiv:2405.04094v1, May 7, 2024. \url{https://arxiv.org/abs/2405.04094v1}
\item
  Victor Y. Wang and Max Wenqiang Xu. \emph{Harper's beyond square-root conjecture}. arXiv:2405.04094v2, final version, August 26, 2025. \url{https://arxiv.org/abs/2405.04094v2}
\item
  Adam J. Harper. \emph{The typical size of character and zeta sums is o(sqrt(x))}. arXiv:2301.04390, January 11, 2023. \url{https://arxiv.org/abs/2301.04390}
\item
  Victor Y. Wang and Max Wenqiang Xu. \emph{Paucity phenomena for polynomial products}. arXiv:2211.02908v1, November 5, 2022. \url{https://arxiv.org/abs/2211.02908v1}
\item
  Regis de la Breteche, Victor Y. Wang, and Max Wenqiang Xu. \emph{Random Multiplicative Functions and Making Squares from Polynomial Values}. arXiv:2607.06398v1, July 7, 2026. \url{https://arxiv.org/abs/2607.06398v1}
\item
  Micah Blumberg. \emph{Neural Lace Podcast 1}. April 11, 2017. \url{https://www.youtube.com/watch?v=YrX_68oKuVs}
\item
  Micah Blumberg. \emph{Coherence and decoherence as complementary information operations}. GitHub commit 37da575b\allowbreak{}a33f85ce\allowbreak{}f4239295\allowbreak{}be161066\allowbreak{}d8d39f4c, June 8, 2022. \url{https://github.com/v5ma/selfawarenetworks/blob/37da575ba33f85cef4239295be161066d8d39f4c/a0142z.md}
\item
  Micah Blumberg. \emph{Time Crystals, Time Dilation, Quantum Gradient Time Dilation}. GitHub commit c4e94e9f\allowbreak{}cd049c97\allowbreak{}d47ef6a5\allowbreak{}f2f8c294\allowbreak{}85f5a98d, July 30, 2022. \url{https://github.com/v5ma/selfawarenetworks/blob/c4e94e9fcd049c97d47ef6a5f2f8c29485f5a98d/a0253z.md}
\item
  Micah Blumberg. \emph{A New Unified Field Theory Called Quantum Gradient Time Crystal Dilation}. SVGN, January 29, 2024. \url{https://www.svgn.io/p/a-new-unified-field-theory-called}
\item
  Micah Blumberg. \emph{Super Information Theory}, version 1. Figshare, February 10, 2025. \url{https://figshare.com/articles/journal_contribution/Super_Information_Theory/28379318/1}
\item
  Self Aware Networks Research Institute. \emph{SAN Claim Genealogy, 2011-2022: Earliest Wording to Mature Mechanism}. Source-governance page, audited July 19, 2026. \url{https://v5ma.github.io/wiki/san/sources/san-claim-genealogy-2011-2022.html}
\end{enumerate}

\end{document}
